\VerdierPhysicalPageAnchor{0025}%
où $\mathrm{RH}(X,{\ })$ désigne le foncteur dérivé total à droite du foncteur
``sections globales'' sur $X$ (de même pour $\mathrm{RH}(Y,{\ })$ et
$\mathrm{RH}(X\times Y,{\ })$) et où $\overset{\mathrm L}{\otimes}_A$
désigne le foncteur dérivé total à gauche du foncteur produit tensoriel. (On
suppose que $X$ et $Y$ sont des espaces de dimension cohomologique finie, par
exemple des complexes cellulaires finis.) Les deux membres de
$\text{(Kü)}_{\mathrm{Der}}$ sont des complexes de $A$-modules
déterminés à quasi-isomorphisme près dont les modules de cohomologie sont
respectivement l'aboutissement de la première et de la deuxième suite
spectrale de $\text{(Kü)}_{\mathrm{coh}}^{4}$. La formule
$\text{(Kü)}_{\mathrm{Der}}$ est, bien entendu, encore valable
lorsque $F$ et $G$ sont des complexes de faisceaux sur $X$ et $Y$
(convenablement bornés). Elle est parfaitement maniable pratiquement, et
permet de formuler également les propriétés de commutativité, et
d'associativité lorsqu'il y a plusieurs facteurs.

\VerdierSemanticAnchor{intro-5-base-change}
L'extension des scalaires conduit à une formule analogue dans les catégories
dérivées. Lorsque $A\to B$ est une $A$-algèbre et lorsque $F$ est un faisceau
de $A$-modules sur $X$, on a la relation :
\[
\mathrm{RH}(X,F)\overset{\mathrm L}{\otimes}_A B
\xrightarrow{\sim}
\mathrm{RH}\bigl(X,F\overset{\mathrm L}{\otimes}_A B\bigr).
\]

\VerdierSemanticAnchor{intro-5-naive-formulas}
De façon imagée, on peut dire, qu'en général,
\emph{les formules naïves et traditionnellement fausses deviennent vraies
lorsqu'on travaille dans les catégories dérivées}.

\VerdierSemanticAnchor{intro-5-l-adic}
L'importance des anneaux de coefficients $\mathbb Z/\ell^r\mathbb Z$ provient
notamment de la cohomologie étale des schémas. On sait [6] que dans cette
théorie on doit essentiellement travailler avec des faisceaux de torsion. On
obtient la cohomologie $\ell$-adique, \emph{i.e.} à valeurs dans l'anneau
$\widehat{\mathbb Z}_{\ell}$ des entiers $\ell$-adiques, en passant à la
limite projective dans la cohomologie à valeurs dans les anneaux
$\mathbb Z/\ell^r\mathbb Z$. La formule
$\text{(Kü)}_{\mathrm{Der}}$ est encore valable dans ce contexte
(lorsque, par exemple, $X$ et $Y$ sont des schémas propres sur un corps
algébriquement clos, $A=\mathbb Z/\ell^r\mathbb Z$), et passe à la limite
immédiatement pour fournir une formule analogue pour les
$\widehat{\mathbb Z}_{\ell}$-faisceaux. Mais comme
$\widehat{\mathbb Z}_{\ell}$ est un anneau principal, la formule
$\text{(Kü)}_{\mathrm{Der}}$ fournit en passant à la cohomologie une
formule du type $\text{(Kü)}_{\mathrm{coh}}^{3}$ et en tensorisant
par $\mathbb Q_{\ell}$, la formule
$\text{(Kü)}_{\mathrm{coh}}^{1}$ pour les
$\mathbb Q_{\ell}$-faisceaux.

\VerdierSemanticAnchor{intro-6}
\noindent\textbf{6.}\quad
C'est le développement systématique des considérations esquissées ci-dessus
que nous présentons dans ce travail. Voici quel en est le plan.

\VerdierSemanticAnchor{intro-6-plan-chapter-one}
Nous étudions au premier chapitre les catégories de complexes des catégories
additives. Ces catégories sont graduées. Plus précisément, elles sont munies
de foncteurs translation (translation des degrés dans une direction
